\phantomsection\bheading{section}{Appendix A. The prior art, and exactly what is taken from each}{section}{Appendix A. The prior art, and exactly what is taken from each}

The reading order of \code{references/\allowbreak{}ownership/\allowbreak{}README.\allowbreak{}md} was followed; each entry says what the design takes, what it declines, and why.

\textbf{Clean} (Barendsen \& Smetsers 1993/1995/1996; Smetsers et al. 1994; de Vries et al. 2006/2007; Clean 2.2 report ch. 9). \emph{Taken:} the demand lives in library types (“newArray :: Int\,$\xrightarrow{\times}$\,Array$^{u}$ rather than Int\,$\xrightarrow{\times}$\,Array$^{\bullet}$”, de Vries 2007 §6, p. 10), so user code carries nothing; the propagation rule for containers — “if a unique object is stored in a data structure, the data structure itself becomes unique as well” (report §9.2, p. 93) — appears here as the holder relation through paths; evaluation-order-aware marking, Smetsers' preliminaries and alternates (1994 §4 Def. 4.2–4.8, p. 9) and the report's “observing” references (§9.4, p. 96), is what §2.3's liveness generalises. \emph{Declined:} attributes on types solved by unification with a subtyping restriction on arrows (1996 Def. 7.6, p. 19), which is why Clean has no principal uniqueness type (“there is no 'Principal Uniqueness Type Theorem'”, p. 32) and must force recursive groups non-unique (de Vries 2007 §7, p. 12); the rank-2 escape for \code{if isEmpty arr then shrink arr else grow arr} (de Vries 2007 §6, p. 11–12), which an annotation-free checker cannot use and liveness makes unnecessary. The \code{foldr} and \code{mark} false errors (Barendsen \& Smetsers 1995 §6, pp. 12–14) are accepted here: \code{foldr}'s callback is $\mathsf{borrowed}$-captured, and \code{mark}'s \code{lookup} precedes its \code{update}.

\textbf{Aspinall, Hofmann, Konečný 2002/2008.} \emph{Taken:} the three aspects (§1, p. 4) as $\mathit{use}$; the LET rule's per-variable condition that a common variable may be read with aspect 2 or 3 in the binding and destroyed in the body (“the modification may happen before the reference”, 2008 Fig. 8, p. 16) — the reads-before-writes rule in its earliest typed form; the best-typing claim and its optimistic search (§6.2, p. 36); the separation theorem's clause C2 (“heap occupied by arguments with aspect 2 or 3 is not modified”, Thm 5.4, p. 23) as Lemma 3's clause 1–2; the two-semantics soundness theorem (Thm 5.8, p. 30) as §4.1's shape. \emph{Declined:} first-order, monomorphic, affine with explicit \code{◇} resources; higher-order deferred to “the standard types and effects technique” (p. 35–36), which §3.6 is.

\textbf{Wadler 1990/1991; Boyland 2001; Marshall, Vollmer \& Orchard 2022.} \emph{Taken:} \code{let!}'s hyperstrict ordering as the reason strictness matters (1990 §4, p. 14); “only update operations are sequentialised but lookups can occur in parallel” (1991 §3, p. 8) as the design goal; alias burying's “all aliases must be dead” (2001 §3.2, p. 13) as the definition of unique-at-a-site; uniqueness as a statement about the past (2022 §2.2, p. 353), which is why a borrowed value \emph{can} be returned to uniqueness after the borrow ends and why \code{List.copy} is the only way back once shared. \emph{Declined:} Boyland's annotations on fields and interfaces (inferred here); Marshall's call-by-name calculus.

\textbf{Wand \& Clinger 2001; Sisal 1990.} \emph{Taken:} liveness-defined destructive update (“x can't be bound to a live location after E1 and E2 have been evaluated”, Def. 8, p. 326) and inclusion constraints with a least solution (“data flow inequalities, which can be solved in polynomial time”, §1.1, p. 320); Sisal's Table I as the model for §9.4's hand classification (25 real sharing, 4 undecidable, of 293). \emph{Declined:} Wand \& Clinger's first-order, scalar-array restriction (§12, p. 344).

\textbf{Lean (Ullrich \& de Moura 2019), Perceus, frame-limited reuse, FP², Brandon et al. 2026.} \emph{Taken:} the direction of inference (“starting with the approximation $\beta(c) = B^n$ … repeat … until we reach a fix point”, 2019 §5.2, p. 6; Brandon's Knaster–Tarski least fixpoint, §4.4, p. 16); the one correctness rule every paper shares — a parameter that reaches a write must be owned (2019 p. 6; Brandon §3.2, p. 7); Brandon's tail-call refinement (“treat variable occurrences in the argument position of a tail-call as escaping”, §6.3, p. 19) which beni's loop rebinding handles directly; FP²'s diagnosis of static uniqueness typing — “code duplication, where a single function can have multiple different implementations” (2023 §1.4.1) — answered here by summaries that are facts about flows, not types, so one \code{reverse} serves both callers and the \emph{caller} is what is checked. \emph{Declined:} run-time reference counts and the silent copy at a shared call (“The fip annotation in Koka only guarantees that no (de)allocation occurs if the parameters are unique at runtime”, Lorenzen, Leijen, Swierstra \& Lindley, PLDI 2024 §3.1, p. 168:7); borrowing as an RC-token discipline that lets a borrowed value be returned with an \code{inc} (2019 Fig. 5); Lorenzen \& Leijen's space objection to inferred borrowing (2022 §4.2) does not apply, because there is no count to keep a borrowed value alive past its last read. Brandon's “a handful of reference count operations” fallback is this design's error.

\textbf{OxCaml (Lorenzen et al. 2024; Peters et al. 2026).} \emph{Taken:} the lock rule for closures — $\mathit{once}^{\dagger}\mathrel{:=}\mathsf{unique},\allowbreak{}\mathsf{many}^{\dagger}\mathrel{:=}\mathsf{aliased}$, a closure's captures are $u_{1}\vee a_{2}^{\dagger}$ (2024 §3.3, p. 10) — as §5.1's $\mathit{once}$ rule; polarised inequality constraints solved by “essentially just transitive closure” (App. B, p. 48) as the shape of §3.6's edges; mode crossing (2026 §2.3, p. 6; §4.2, p. 18) as §2.1's $\mathit{crosses}(\tau )$; the measured annotation counts (27 382 locality annotations in interfaces; “the implementations can generally infer”, 2024 §7, p. 24) as the warning about interface churn; the currying remark (§6.4, p. 22) as evidence that beni's no-currying decision removes the hard case. \emph{Declined:} modes as types with no polymorphism (“85 functions duplicated”, §6.6, p. 23) — summaries here are per flow, not per mode; the per-module inference with written interface modes (docs, \emph{pitfalls}) — summaries here are inferred \emph{and} published.

\textbf{Reachability types (Bao 2021; Wei 2024; Jia 2026; O'Connor 2026; Capybara; Deng 2025).} \emph{Taken:} a function type that names which argument its result reaches ($f(x:T^{\blacklozenge})\to T^{x}$, Wei 2024 §2.2.3) as $\mathit{res}$; the warning that polymorphism over tracked/untracked values is unsound without a fresh marker (\code{fakeid}, ms.tex:1179–1216) — here $\mathsf{fresh}$ is a distinct cell kind, never a type variable; “a consume kills its aliases” (Capybara ms.tex:2516) as the dead-alias lemma's content; “a component of a shared structure cannot be consumed directly” (Capybara §7) as §5.4's rule for containers; Deng's use/kill effects on closures (Fig. 6, p. 11) as the $\mathit{cap}$ facts, inferred rather than declared; O'Connor's statement of the guarantee. \emph{Declined:} qualifier annotations on arguments and explicit instantiation (Jia 2026 §4: inference is given “annotations on function arguments and explicit instantiations”); the opaque joint qualifier for lists ($\mu h.\mathit{List}[\mathit{Ref}^{h}]$), replaced by per-path cells.

\textbf{Rust: NLL, Polonius, two-phase borrows, Oxide, Pearce, Ho, Emre.} \emph{Taken:} liveness as the precision lever (“if a variable is live on entry to a point P, then all regions in its type must include P”, 2017 post) and the minimal-lifetime principle (RFC 2094); two-phase borrows as the formal reads-before-writes (“acts exactly like a shared borrow” until activation, RFC 2025); the three-point error with the blame on the action (RFC 2094, \emph{Leveraging intuition}); Polonius as “a bunch of transitive closures” over directional facts (Stjerna §3.3, p. 15); Pearce's and Ho's joins at merges as the model for §5.6; Emre's measurement as the argument against unification (§4, Table 2). \emph{Declined:} lifetimes in signatures; references into structures (beni has none, which is why problem case \#3 does not arise).

\textbf{Futhark.} \emph{Taken:} the consumption judgement $(O_{2}\cup C_{2})\cap C_{1}=\emptyset$ (PLDI 2017 Fig. 6, p. 7) as the one-line statement of the rule; \code{if} as a union of alias sets (ALIAS-IF) for cells, with per-path liveness added; the refusal of a free-variable consume inside a lambda (§3.3, p. 6) as the $\mathit{once}$ rule; the separate pass after type checking (2026 rewrite post); the error index's shape — message, minimal program, reason, \code{copy} fix, and when the copy “subverts the purpose” (§8.1.2); the 2021 lesson on naming intermediate results; the free-theorem idea for \code{id} and \code{|>}. \emph{Declined:} \code{*} on parameters and results, intra-procedural analysis with no summaries (“we are forced to be conservative”, §3.2, p. 6), the 2026 verdict “don't do it” — which was about a type system reasoning about aliasing, and §7.1's point is that this is not a type system.

\textbf{Usage analysis (Turner 1995; Wansbrough \& Peyton Jones 1999/2002); Affe (Radanne 2020); escape analysis (Blanchet 1998/1999).} \emph{Taken:} the poisoning problem as the reason summaries must be polymorphic over function arguments (“one call to f 'poisons' all the others”, 1999 §6.3); polarity simplification of constraint sets (thesis §4.1.3, p. 95) as the way to keep summary blocks small; the measured warning that generalising local binders loses precision; Affe's \code{Abs} rule where a captured shared borrow leaves the closure unrestricted (§3.3) as the shape of $\mathit{cap}=\mathsf{borrowed}$; Blanchet's closure case split — applied versus escaped — as $\mathit{escapes}$. \emph{Declined:} usage on every type (demand-counting, which counts a read-only capture as a use); Affe's lexical regions and explicit \code{\&x}.

\textbf{Error design (Czaplicki 2015; Wrenn \& Krishnamurthi 2017; Marceau 2011; rustc guide; Crichton 2020/2023/2024; Hylo; Swift).} Taken in full in §6; the one dissent recorded is Marceau's “error messages should not propose solutions” for novices, against which Hylo, Futhark, Rust and Swift all ship fixes; this document ships the fix and the verdict line.

\textbf{Research 42 and 68/69.} Research 42's R0 is this checker's backend half: “in place where last use and per-function interface summaries prove a value unique” (§1), with S2's “consumed-unique parameters … summaries go into the interface and flow forward” and S4's result summaries; its findings that the TEA model needs a hands-over runtime (§6.1) and that a copy at an unprovable call is “catastrophic in a loop” (§6.4) are why the rule is an error at the site and not a copy. Research 69 §1.3's nine items are each answered: 1 (§2.6), 2 (§2.3), 3 (§2.4), 4 (§3.5), 5 (§3.9), 6 (§6.5), 7 (§2.1), 8 (§7.1), 9 (§5.1).
